\documentclass[11pt]{article}
\usepackage[margin=1in]{geometry}
\usepackage[T1]{fontenc}
\usepackage{lmodern}
\usepackage{amsmath,amssymb,amsthm}
\usepackage[hidelinks]{hyperref}
\newtheorem{theorem}{Theorem}
\newtheorem{lemma}{Lemma}
\newtheorem{corollary}{Corollary}
\title{NP-completeness of aggregate root-dependent spanning-tree cost}
\author{Alec Kriebel\\Independent researcher\\
\small\href{https://orcid.org/0009-0001-9320-500X}{ORCID: 0009-0001-9320-500X}}
\date{6 October 2026}
\begin{document}
\maketitle
\begin{abstract}
An undirected spanning tree induces an arborescence at every vertex root.
Given an independent arc-cost vector for each root, we prove that deciding
whether the sum of the costs of these entire rooted orientations is at most
a given threshold is NP-complete, even when every cost belongs to
$\{0,1,2\}$. A reduction from satisfiability uses two hubs, variable vertices,
and clause vertices. A symmetric cost vector at one hub forces a tree at
threshold to encode an assignment with one selected literal per clause;
the corresponding clause-root costs test those literals. Uniformly shifting
the costs yields the same conclusion for costs in $\{1,2,3\}$. This resolves
Problem~1 in Volker Kaibel's November 2018 Oberwolfach workshop contribution,
whose report was published in December 2019. Exact finite checks accompany
the all-size proof.
\end{abstract}

\section{Problem and result}
Let $G=(V,E)$ be a finite connected simple undirected graph with
$N=|V|\geq 1$, and let
$A=\{(u,v),(v,u):\{u,v\}\in E\}$. For every $r\in V$, the input supplies
an independent cost vector $c_r:A\to\mathbb Z_{\geq0}$. All
$2N|E|$ entries are explicitly encoded in binary. For a spanning tree $T$,
let $T^r$ denote its entire edge set oriented away from $r$, and define
\begin{equation}\label{eq:objective}
 F(T)=\sum_{r\in V}\ \sum_{(u,v)\in T^r}c_r(u,v).
\end{equation}
The decision question is whether $F(T)\leq K$ for some spanning tree, where
$K\in\mathbb Z$ is part of the input. The same common undirected tree
is used at every root; the costs charge all its oriented edges.
Reversing every input cost, $c'_r(u,v)=c_r(v,u)$, gives the equivalent
inward-orientation convention.

Kaibel's Problem~1 asks for the hardness of minimizing precisely this
all-root objective with real arc-cost vectors~\cite{kaibel}. The
nonnegative integral subclass below therefore answers that optimization
question under ordinary finite input encoding. Kaibel motivates the question
by integer optimization over Martin's extension of the spanning-tree
polytope. That formulation is classical, as described in the inspected
expositions of Conforti, Cornu\'ejols and Zambelli~\cite[Section~6.1]{ccz}
and Friesen~\cite[Proposition~2]{friesen}: common edge variables satisfy
$z_r(u,v)+z_r(v,u)=x_{\{u,v\}}$, with $z\geq0$, root indegree zero and all
other indegrees one. For integral common edge variables these constraints select
one tree, and its root orientations are determined uniquely. Integrality
of the projection onto the edge variables does not establish integrality
of the full joint lift under arbitrary objectives on all $z$ variables.
The formulation and this interpretation are not contributions of this note.

\begin{theorem}\label{thm:main}
The decision problem for~\eqref{eq:objective} is NP-complete even when
every cost belongs to $\{0,1,2\}$. Consequently it is strongly NP-complete,
and finding a minimum-cost tree is NP-hard.
\end{theorem}

Membership in NP follows by checking that the proposed $N-1$ edges form
a spanning tree, traversing it from each root, and summing the supplied
binary costs. The number of operations and the lengths of intermediate
integers are polynomial in the explicit input size.

\section{Reduction}
We reduce from satisfiability of a conjunctive normal form whose clauses
have at most three literals, using the standard NP-completeness result
for this restriction~\cite{karp}. Remove repeated literals within each
clause and discard tautological clauses. An empty clause gives a fixed
no-instance: two vertices joined by one edge, all four root-arc costs
equal to $1$, and threshold $1$. Its unique tree has aggregate cost $2$.
If no clauses remain, use the same graph with all costs zero and threshold
zero, a fixed yes-instance. Delete unused variable declarations and
relabel the variables that actually occur densely. Thus $n$ is the number
of occurring variables, regardless of their original labels. It remains
to consider $n\geq1$ variables $x_1,\ldots,x_n$ and $m\geq1$ nonempty
clauses $C_1,\ldots,C_m$, each with at most three distinct variables and
without both signs of a variable in one clause.

The graph has two hubs $t,f$, a vertex $v_i$ for every variable, and a
vertex $q_j$ for every clause. Its edges are
\[
 \{t,f\},\qquad \{t,v_i\},\{f,v_i\}\ (1\leq i\leq n),\qquad
 \{q_j,v_i\}\ \text{when a literal on }x_i\text{ occurs in }C_j.
\]
It is connected and simple, with $N=n+m+2$ vertices and at most
$1+2n+3m$ edges. Set
\begin{equation}\label{eq:threshold}
 K=2m+n.
\end{equation}

At root $t$, assign the same cost to both orientations of each edge:
\begin{equation}\label{eq:basecost}
 \begin{aligned}
 c_t(t,f)=c_t(f,t)&=0,\\
 c_t(t,v_i)=c_t(v_i,t)=c_t(f,v_i)=c_t(v_i,f)&=1,\\
 c_t(q_j,v_i)=c_t(v_i,q_j)&=2.
 \end{aligned}
\end{equation}
At root $q_j$, the only potentially nonzero costs are variable-to-hub
arcs for variables occurring in $C_j$:
\begin{equation}\label{eq:literalcost}
 \begin{array}{c|cc}
 \text{literal in }C_j & c_{q_j}(v_i,t)&c_{q_j}(v_i,f)\\\hline
 x_i&0&1\\
 \neg x_i&1&0
 \end{array}
\end{equation}
All other costs at $q_j$ are zero, including the reverse hub-to-variable
arcs. The vectors at $f$ and all variable roots are identically zero.
These prescriptions specify every entry of the input table, without
using any tree-dependent data.

\begin{lemma}[Structural threshold]\label{lem:structure}
If a spanning tree has $F(T)\leq K$, it contains the hub edge, every clause
vertex is a leaf, and every variable has exactly one hub attachment.
\end{lemma}
\begin{proof}
Let $q$ be the number of clause-variable edges in $T$, $p$ the number
of hub-variable edges, and $h\in\{0,1\}$ the indicator of the hub edge.
Each clause vertex has positive degree and is incident only to
clause-variable edges, each of which meets exactly one clause vertex.
Hence $q\geq m$. Counting all edges in the tree gives
$p+q+h=n+m+1$. By the symmetric costs at $t$,
\begin{equation}\label{eq:structural}
 F_t(T):=\sum_{a\in T^t}c_t(a)=p+2q
       =K+(1-h)+(q-m).
\end{equation}
All other root costs are nonnegative. The two terms after $K$ are
nonnegative integers, so $F(T)\leq K$ forces $h=1$ and $q=m$.
Every clause vertex is therefore a leaf. Deleting these leaves leaves a
tree on the hubs and variables containing the hub edge. Each variable
must meet at least one hub, and meeting both would create a triangle
with the hub edge. Each meets exactly one.
\end{proof}

Call a tree of this form structured. Encode an assignment by making $x_i$
true exactly when $\{t,v_i\}$ is present. Each clause leaf selects its
unique neighboring variable and hence its unique literal in that clause.
Conversely, any assignment together with one selected literal in each
clause gives a structured tree, with $F_t(T)=K$.

\begin{lemma}[Literal costs]\label{lem:literal}
For every structured tree,
\begin{equation}\label{eq:structured}
 F(T)=K+\#\{j:\text{the literal selected by }q_j\text{ is false}\}.
\end{equation}
\end{lemma}
\begin{proof}
Root the structured tree at $q_j$, and let $v_i$ be its unique neighbor.
The path into the hub core begins $q_j\to v_i\to h_i$, where $h_i$ is
the selected hub of $v_i$. Thus this variable's hub edge is directed
from the variable to its hub. Every other variable's hub edge is
directed from its hub toward that variable: the core consists of two
stars joined by the hub edge, and the route from $q_j$ enters every
other variable from the hub side. Accordingly, the only arc that can
have nonzero cost at root $q_j$ is $v_i\to h_i$.
Table~\eqref{eq:literalcost} assigns this arc cost zero precisely when
the selected literal is true, and one otherwise. All clause edges,
the hub edge, and the remaining root vectors contribute zero beyond
the base cost. Summing over the clause roots proves the identity.
\end{proof}

If the formula is satisfiable, attach every variable to its truth-value
hub and attach every clause leaf to a true literal. The resulting tree
has cost $K$. Conversely, a tree of cost at most $K$ is structured by
Lemma~\ref{lem:structure}; Lemma~\ref{lem:literal} makes each selected
literal true, so every clause is satisfied. This proves the equivalence
of the two decision instances.

The graph, all $N$ cost vectors, and the threshold are computable in
polynomial time. There are $2N|E|$ explicitly supplied entries, each in
$\{0,1,2\}$, and $K=2m+n\leq 2N$. This completes the proof of
Theorem~\ref{thm:main}, including hardness under unary numerical encoding.

\begin{corollary}\label{cor:positive}
The same decision problem is strongly NP-complete when every cost belongs
to $\{1,2,3\}$.
\end{corollary}
\begin{proof}
Add one to every arc cost at every root and replace $K$ by
$K+N(N-1)$. Every spanning tree has $N-1$ arcs at each of the $N$ roots,
so every objective changes by exactly this same constant.
\end{proof}

The original audited choice placed $B=n+1$ on clause-variable edges.
For that same construction at any integer $B\geq2$, with $K_B=Bm+n$,
the unrestricted base identity is
\[
 F_t(T)=p+Bq=K_B+(1-h)+(B-1)(q-m).
\]
The choice $B=2$ above is a consequence of that identity and the same
literal test. The transpose and positive shift are likewise consequences
of the same result. Equation~\eqref{eq:structured} is asserted only for
structured trees; it does not describe the unrestricted global optimum
by the minimum number of unsatisfied clauses.

\section{Verification, attribution and status}
The accompanying standard-library Python package constructs full explicit
cost tables and enumerates all spanning trees for small test instances.
It sums costs by independently orienting each whole tree from every root
and also by partitioning roots across each edge cut. It checks the
unrestricted base identity, individual clause-root costs, both reduction
directions, inward transposition, uniform positive shifts, normalization
with sparse variable labels, and trivial cases. Deliberately false guards,
corrupted costs and an incorrect proof binding are rejected under normal
Python and Python \texttt{-O}. Finite computations supplement the written
all-size proof; they do not establish NP-hardness or historical priority.

The question is attributed to Kaibel's workshop contribution of
4--10 November 2018, printed on pages 3014--3015 of the report published
17 December 2019~\cite{kaibel}. No earlier proof of the exact objective
was located in the source and model audit completed on 6 October 2026.
That dated audit has incomplete historical coverage and does not certify
absolute first discovery or continued open status in 2026. Its source
scopes and access limitations are recorded in the supporting provenance.
The result concerns this exact whole-orientation objective; no claim is
made here about Kaibel's neighboring Problem~2.

\paragraph{AI assistance and review.}
AI tools were used extensively in solving, drafting, literature checking,
code generation, verification and adversarial audits. Automated audits
are not conventional human peer review. This manuscript has not received
conventional human peer review or refereeing.

\begin{thebibliography}{9}
\bibitem{kaibel}
V. Kaibel, Open Problem: Arborescences and NP-hardness, in
\emph{Combinatorial Optimization}, Oberwolfach Reports 15 (2018),
Report 50, pp. 3014--3015. Workshop 4--10 November 2018;
report published 17 December 2019.
\href{https://doi.org/10.4171/OWR/2018/50}{doi:10.4171/OWR/2018/50}.

\bibitem{karp}
R. M. Karp, Reducibility among combinatorial problems, in
R. E. Miller and J. W. Thatcher (eds.), \emph{Complexity of Computer
Computations}, Plenum Press, 1972, pp. 85--103.
\href{https://doi.org/10.1007/978-1-4684-2001-2_9}{doi:10.1007/978-1-4684-2001-2\_9}.

\bibitem{ccz}
M. Conforti, G. Cornu\'ejols and G. Zambelli,
\emph{Extended Formulations in Combinatorial Optimization},
November 2009, revised February 2010, Section 6.1, pp. 25--26.
\href{https://www.andrew.cmu.edu/user/gc0v/webpub/ExtFor-Feb2010.pdf}{Primary author manuscript}.

\bibitem{friesen}
M. Friesen, \emph{Extended Formulations for Higher Order Polytopes in
Combinatorial Optimization}, dissertation, Otto-von-Guericke-Universit\"at
Magdeburg, 2019, Proposition 2, p. 8.
\href{https://d-nb.info/1219966134/34}{Primary dissertation}.
\end{thebibliography}
\end{document}
